% Standalone source. Compile with pdflatex twice; the bibliography is inline.
\documentclass[11pt,a4paper]{article}
\usepackage[T1]{fontenc}
\usepackage{lmodern}
\usepackage[margin=26mm,top=25mm,bottom=26mm]{geometry}
\usepackage{amsmath,amssymb,amsthm,microtype}
\usepackage{xcolor}
\usepackage{fancyhdr}
\usepackage{hyperref}
\definecolor{inklink}{RGB}{24,74,96}
\hypersetup{colorlinks=true,linkcolor=inklink,citecolor=inklink,urlcolor=inklink,
 pdftitle={A proof of the Caporossi-Hansen conjecture on the Randic index and graph radius},
 pdfauthor={Keio and Shaitetsu},pdfkeywords={Randic index, vertex radius, absolute radius, degree cut}}
\pagestyle{fancy}\fancyhf{}\setlength{\headheight}{14pt}
\fancyhead[L]{\small\itshape The Caporossi--Hansen conjecture}
\fancyhead[R]{\small Keio and Shaitetsu}\fancyfoot[C]{\thepage}
\renewcommand{\headrulewidth}{0.25pt}
\newtheorem{theorem}{Theorem}[section]
\newtheorem{lemma}[theorem]{Lemma}
\newtheorem{proposition}[theorem]{Proposition}
\newtheorem{corollary}[theorem]{Corollary}
\newcommand{\rad}{\operatorname{rad}}
\newcommand{\diam}{\operatorname{diam}}
\newcommand{\ecc}{\operatorname{ecc}}
\newcommand{\dist}{\operatorname{dist}}
\allowdisplaybreaks[1]
\setlength{\emergencystretch}{1em}
\title{A proof of the Caporossi--Hansen conjecture\\on the Randi\'c index and graph radius}
\author{Keio \and Shaitetsu}
\date{13 November 2026\quad$\cdot$\quad Research manuscript for independent review}
\begin{document}
\maketitle
\begin{abstract}
For every finite connected simple graph of maximum degree $\Delta\ge2$,
we prove $R(G)-\rad(G)\ge\sqrt{\Delta}-3/2$, where $R$ is the classical
Randi\'c index. Thus the paths of even order at least four are exactly
the connected graphs whose Randi\'c index is smaller than their vertex
radius. Equality between the index and vertex radius holds exactly for the singleton, the single edge and
the even cycles. The main estimate bounds the absolute radius $\rho(G)$
by the sum $M(G)$ of reciprocal maximum endpoint degrees. In a minimal counterexample,
deleting the leaves attached to vertices below a degree threshold gives
an upper bound for the number of edges crossing that threshold. Summing
these bounds and constructing a spanning tree of the two-connected block
gives $2\rho(G)\le n-\Delta+4\le2M(G)$ in the proposed counterexample,
where $n$ is the order of that block.
\end{abstract}

\section{Introduction and principal results}
For a finite connected undirected simple graph $G$, let $d_v$ denote the
degree of a vertex $v$. Its Randi\'c index and its vertex radius are
\[
 R(G)=\sum_{uv\in E(G)}\frac{1}{\sqrt{d_ud_v}},\qquad
 \rad(G)=\min_{v\in V(G)}\max_{w\in V(G)}\dist_G(v,w).
\]
Here $uv$ denotes an unordered edge, counted once. The one-vertex graph
has radius zero and an empty edge sum. The notation $P_N$ and $C_N$
denotes a path and a cycle on $N$ vertices, respectively. Throughout,
$\Delta$ denotes maximum degree and $\diam$ denotes diameter. The usual
bound $\diam(G)\le2\rad(G)$ follows by joining two vertices through a
vertex center.

Fajtlowicz proposed $R(G)\ge\rad(G)-1$ \cite{Faj}. Caporossi and Hansen
conjectured $R(G)\ge\rad(G)$ for every connected graph other than a
path of even order and proved the tree case \cite{CH,Zuo}. We obtain a
bound in terms of maximum degree, as well as a complete classification.
\begin{theorem}\label{thm:main}
Let $G$ be a finite connected simple graph with maximum degree $\Delta\ge2$.
Then
\begin{equation}\label{eq:main}
 R(G)-\rad(G)\ \ge\ \sqrt{\Delta}-\tfrac32.
\end{equation}
For the class of all finite connected simple graphs, the following alternatives hold.
\begin{enumerate}
\item $R(G)<\rad(G)$ precisely for $P_{2m}$ with $m\ge2$.
\item $R(G)=\rad(G)$ precisely for $K_1$, $K_2$, and $C_{2m}$ with $m\ge2$.
\item Every other graph satisfies $R(G)>\rad(G)$.
\end{enumerate}
The minimum of $R(G)-\rad(G)$ over all such graphs is $\sqrt2-3/2$,
attained precisely in the first case.
\end{theorem}

The modified Randi\'c index
\[
 M(G)=\sum_{uv\in E(G)}\frac1{\max(d_u,d_v)}
\]
occurs in the diameter analysis of Dvo\v{r}\'ak, Lidick\'y and
\v{S}krekovski \cite{DLS} and in the radius argument of Cygan, Pilipczuk
and \v{S}krekovski \cite{CPS}. The latter proved the Caporossi--Hansen
inequality for graphs of maximum degree at most four and established
$M(G)\ge\rad(G)-1/2$ in that class. Zuo proved the conjecture for
unicyclic and bicyclic graphs \cite{Zuo}. Deng, Tang and Zhang gave a
bound depending on the cyclomatic number \cite{Deng}. Doig studied the
inequality for cactus graphs \cite{Doig}. The spanning-tree estimate used
in Section~\ref{sec:metric} also appears as Lemma~2.1 of \cite{CPS}.

Give every edge of $G$ length one and denote the resulting metric graph
by $|G|$. Write
\[
 \rho(G)=\min_{x\in|G|}\ \max_{w\in V(G)}\dist_{|G|}(x,w).
\]
Thus a center may lie inside an edge, and the vertices are the demand points.
\begin{theorem}[Absolute-radius inequality]\label{thm:abs}
Every finite connected simple graph satisfies $M(G)\ge\rho(G)$.
\end{theorem}
The proof of Theorem~\ref{thm:abs} builds on the minimal-counterexample
arguments in VMRC [087] and [106] and the local degree identity in
VMRC [094] \cite{VMRC}. All reductions are proved in
Sections~\ref{sec:metric}--\ref{sec:sum}. The core with attached leaves
is a bag in the sense of Definition~4.1 of \cite{CPS}. The step that
removes the degree
bound is Proposition~\ref{prop:cut}, which compares each cut with the effect
of deleting just the leaves on its lower-degree side. Section~\ref{sec:finish}
then deduces Theorem~\ref{thm:main} by a comparison edge by edge.

\section{Metric centers and spanning trees}\label{sec:metric}
Distances between vertices in $|G|$ agree with graph distances. With
finitely many vertices, the maximum of their distance functions is a
continuous function on the compact space $|G|$, so an absolute center
exists.
\begin{lemma}\label{lem:half}
An absolute center can be chosen at a vertex or an edge midpoint.
In particular $\rad(G)-1/2\le\rho(G)\le\rad(G)$.
\end{lemma}
\begin{proof}
Consider an edge $uv$ with coordinate $t\in[0,1]$ measured from $u$.
For each vertex $w$ the distance from $t$ to $w$ is
\[
 \min\{\dist_G(u,w)+t,\ \dist_G(v,w)+1-t\}.
\]
Both affine functions have integer intercepts and slopes $\pm1$. Two
such functions of opposite slopes can meet only at a half-integer.
Consequently every individual distance, and their maximum, is affine
on each of the open half-edges. Its minimum on the closed edge is
attained at an endpoint or midpoint. Distances from a center so chosen
to vertices are multiples of $1/2$. Rounding any absolute center to a
nearest vertex increases all its distances by at most $1/2$. The vertex
centers are allowed when minimizing for $\rho$.
\end{proof}
\begin{lemma}\label{lem:tree}
For a tree $T$ of diameter $L$, one has $\rho(T)=L/2$ and
$\rad(T)=\lceil L/2\rceil$. If a connected graph $H$ on $m$ vertices
has a vertex with at least $j\ge2$ incident edges, then it admits a
spanning tree of diameter at most $m-j+1$. In particular
$2\rad(H)\le m-j+2$.
\end{lemma}
\begin{proof}
Let a longest path in $T$ have endpoints $a,b$. For any vertex $w$, its
branch meets the path at a unique point at distance $s$ from $a$; write
$t$ for the distance from that point to $w$. The two paths to the
endpoints give $t\le\min(s,L-s)$. Thus the distance from the midpoint
to $w$ is at most $|s-L/2|+\min(s,L-s)=L/2$. The endpoints force a
radius of at least $L/2$ by the triangle inequality. Rounding the
midpoint yields the asserted integer radius.

The $j$ chosen incident edges form a tree. Extend this tree, one edge
to a new vertex at a time, to a spanning tree $T$ of $H$. At least $j$
components of $T-v$ meet the chosen vertex $v$ by distinct edges.
Each contains a leaf of $T$, as a farthest vertex of that component
from $v$ is a leaf. A longest path in $T$ uses at most two leaves. It
therefore uses at most $m-(j-2)$ vertices and has at most $m-j+1$ edges.
Distances in $H$ can only be shorter than those in its spanning tree.
The vertex-radius bound now follows from $2\lceil L/2\rceil\le L+1$.
\end{proof}

\section{Reductions of a minimal counterexample}\label{sec:red}
Suppose Theorem~\ref{thm:abs} fails. Among counterexamples $G$ choose
one minimizing the pair $(\beta(G),|E(G)|)$ lexicographically, where
$\beta(G)=|E(G)|-|V(G)|+1$ is the cyclomatic number. This order on
pairs of nonnegative integers is well founded. A graph smaller in this
order satisfies $M\ge\rho$. Stars with at least two leaves have
$M=\rho=1$. On a path of order $m\ge3$, each edge contributes $1/2$
to $M$ and Lemma~\ref{lem:tree} gives $M=\rho=(m-1)/2$. Also
$(M,\rho)=(0,0)$ for $K_1$ and $(1,1/2)$ for $K_2$.

\begin{lemma}\label{lem:delete}
If $e$ is a nonbridge edge of $G$, then $M(G-e)>M(G)$. If $z$ is
a leaf with neighbor $v$ of degree $d\ge2$, and $b$ of the remaining
neighbors of $v$ have degree less than $d$, then
\begin{equation}\label{eq:oneleaf}
 M(G)-M(G-z)=\frac1d-\frac{b}{d(d-1)}\ge0,\qquad
 0\le\rho(G)-\rho(G-z)\le\tfrac12.
\end{equation}
The leaf formulas apply to any connected graph in the stated range.
\end{lemma}
\begin{proof}
Deleting a nonbridge preserves the demand vertices, removes potential
center locations, and cannot shorten distances. It reduces $\beta$ by
one. Hence $M(G-e)\ge\rho(G-e)\ge\rho(G)>M(G)$.

Deleting $z$ removes $1/d$ from $M$. If the other endpoint of an edge
at $v$ has degree less than $d$, its weight increases from $1/d$ to
$1/(d-1)$. The other weights stay fixed. As $b\le d-1$, this proves
the index formula.

Project a center lying on the removed pendant interval to $v$. All
distances to retained vertices weakly decrease, proving the first
radius inequality. Conversely, choose a vertex or midpoint center $y$
of $G-z$, and put $r_0=\rho(G-z)$. If $\dist(y,v)\le r_0-1/2$,
then $\dist(y,z)\le r_0+1/2$ after restoration. Otherwise
$\dist(y,v)=r_0$ by Lemma~\ref{lem:half}. As $G-z$ has at least two
vertices, $r_0\ge1/2$. Move $y$ a distance $1/2$ along a shortest
path toward $v$. The distances to retained vertices increase by at
most $1/2$, and the distance to $z$ is $r_0+1/2$.
\end{proof}

A vertex adjacent to a leaf will be called its support. A support of
degree two in $G$ has a second neighbor of degree at least two, because
$G$ is not a star. Formula~\eqref{eq:oneleaf} then gives a drop in
$M$ of $1/2$, at least the drop in $\rho$. Deleting this leaf would
produce a smaller counterexample with the same $\beta$. Thus every
support in $G$ has degree at least three.

If $z,z'$ are leaves with the same support $v$, delete $z$. In the
remaining graph a center can be chosen outside the pendant interval
leading to $z'$, except possibly at $v$. Indeed, projecting a point
of that interval to $v$ decreases distances to all other vertices,
and puts $z'$ at distance one. Any point on that interval has
eccentricity at least one, since another neighbor of $v$ remains.
Restoring $z$ now gives $z$ the same distance from the center as $z'$.
The radius is unchanged and $M$ weakly decreases upon deletion.
Minimality gives a contradiction. Each support therefore has exactly
one leaf.

\begin{lemma}[Splitting at a cut vertex]\label{lem:split}
For every vertex $v$ of the minimal counterexample $G$, at most one
component of $G-v$ contains an edge.
\end{lemma}
\begin{proof}
Otherwise partition $V(G)\setminus\{v\}$ into unions $A,B$ of components
of $G-v$, each containing an edge. For $i\in\{A,B\}$ let $t_i$ be
the number of neighbors of $v$ in $i$, and $m_i$ the number of edges
within $i$. All four numbers are positive. Form $G_A$ by retaining
$G[A\cup\{v\}]$ and adding $t_B$ new leaves at $v$. Define $G_B$
symmetrically. These graphs are connected, $|E(G_A)|=|E(G)|-m_B$,
and $\beta(G_A)=m_A+t_A-|A|$. The analogous expression for $G_B$
gives $\beta(G_A)+\beta(G_B)=\beta(G)$. Both graphs are therefore
smaller. Each retained edge has its original endpoint degrees.
The $t_A+t_B=d_v$ artificial edges across the two graphs each have
weight $1/d_v$, which gives
\begin{equation}\label{eq:splitindex}
 M(G_A)+M(G_B)=M(G)+1 .
\end{equation}

Set $r_i=\rho(G_i)$. In each $G_i$, an artificial leaf and a retained
neighbor of $v$ are at distance two, so $r_i\ge1$. Choose a center
$y_i$ in the retained part: if a center lies on an artificial pendant
edge, project it to $v$. This weakly decreases its distances to all
vertices other than the tip of that edge, and places that tip at
distance $1\le r_i$. If $s_i=\dist(y_i,v)$, the distance to an
artificial leaf is $s_i+1$, so $s_i\le r_i-1$.

In $|G|$ join $y_A$ to $v$ and $v$ to $y_B$ by shortest paths. Their
concatenation is a geodesic since the two retained parts meet only at
$v$. With signed distance $u$ from $v$ along this path, its ends have
coordinates $-s_A,s_B$. Via $y_A$ and $y_B$, respectively, the
distances to the original vertices on the two sides (including $v$)
are at most
\[
 L(u)=r_A+s_A+u,\qquad H(u)=r_B+s_B-u .
\]
If these functions meet on the interval, their common value is
$(r_A+s_A+r_B+s_B)/2\le r_A+r_B-1$. If $L$ dominates throughout,
at $u=-s_A$ their maximum is $r_A$. If $H$ dominates throughout,
at $u=s_B$ the maximum is $r_B$. Both are at most $r_A+r_B-1$.
Consequently $\rho(G)\le r_A+r_B-1\le M(G_A)+M(G_B)-1=M(G)$ by
minimality and \eqref{eq:splitindex}, a contradiction.
\end{proof}

\begin{proposition}\label{prop:core}
The minimal counterexample consists of a two-connected induced subgraph
$K$ of order at least three together with at most one leaf attached to
each vertex of $K$. Every edge of $K$ is a nonbridge in $G$.
\end{proposition}
\begin{proof}
A tree containing a path of length four violates Lemma~\ref{lem:split}
at the middle of that path. A tree of diameter at most two is a
singleton, an edge or a star. In a tree of diameter three, every vertex
lies on the central edge of a diametral path or is adjacent to one of
its ends: the midpoint covers the tree within $3/2$ by
Lemma~\ref{lem:tree}. Each central vertex supports at most one leaf,
so this tree is a path. These possibilities were excluded above.

Thus $G$ contains a cycle. Its vertices induce a two-connected graph,
where two-connected means having at least three vertices and remaining
connected after deletion of any one vertex. Extend that set to a
maximal vertex set inducing a two-connected subgraph $K$. If a component
outside $K$ has two distinct attachment vertices in $K$, it gives a
path between them with all interior vertices outside $K$. Adding the
path preserves two-connectivity: deletion of an interior vertex leaves
each remaining piece attached at an end, and deletion of a vertex of
$K$ leaves $K$ connected with at least one end of the path attached.
This contradicts maximality. Each outside component therefore has a
unique attachment $v$. If it contained an edge, deleting $v$ would
separate that edge from $K-v$, which is connected and has at least
two vertices. Lemma~\ref{lem:split} excludes this. Each outside
component is a single leaf, and the previous reduction bounds their
number per support. Finally, every vertex of $K$ has at least two
neighbors in $K$. For a core edge $uv$, choose another core neighbor
$w$ of $u$; a path in $K-u$ from $w$ to $v$ completes an alternative
route from $u$ to $v$. Hence $uv$ is a nonbridge.
\end{proof}

\section{A bound for each degree cut}\label{sec:cuts}
Fix the subgraph $K$ of Proposition~\ref{prop:core} and write
$n=|V(K)|$. Degrees $d_v$ in this section are always measured in $G$.
At a core vertex $v$, let $p_v$ be $1$ if $v$ supports a leaf and
$0$ otherwise. Define
\[
 \ell_v=\#\{w\sim v:d_w<d_v\},\quad h_v=d_v-\ell_v,\quad
 g_v=d_v(d_v-1),\quad a_v=\ell_v/g_v,\quad b_v=h_v/g_v .
\]
In particular $\ell_v$ includes the leaf when $p_v=1$.

The next lemma includes the local-minimum and minimum-degree properties
of bags from Lemma~4.2 of \cite{CPS}.
\begin{lemma}\label{lem:local}
Every core vertex has $d_v\ge3$ and $\ell_v\ge1$. For a core edge $uv$
with $d_u<d_v$ we have $a_u>b_v$.
\end{lemma}
\begin{proof}
On deleting $uv$, each remaining edge at an endpoint $x$ gains $1/g_x$
precisely when its other endpoint has degree strictly below $d_x$.
No remaining edge occurs at both $u$ and $v$, since $G$ is simple.
Subtracting the weight of the deleted edge gives
\begin{align}
 M(G-uv)-M(G)
 &=a_u+\frac{\ell_v-1}{g_v}-\frac1{d_v}
   =a_u-b_v &&(d_u<d_v),\label{eq:strict}\\
 &=\frac{\ell_u+\ell_v}{d(d-1)}-\frac1d
   &&(d_u=d_v=d).\label{eq:equal}
\end{align}
Both changes must be positive by Lemma~\ref{lem:delete}. If
$\ell_u=0$, there is no leaf at $u$. A core neighbor of higher degree
would contradict \eqref{eq:strict}, since $b_v\ge0$. A neighbor of
equal degree would have $\ell_v\le d-1$, contradicting
\eqref{eq:equal}. These exhaust the neighbors of $u$. Finally, a
vertex of full degree two already has two neighbors in $K$; all
vertices of $K$ have degree at least two. It could have no lower-degree
neighbor.
\end{proof}

A core vertex of minimum full degree must have a leaf, by
Lemma~\ref{lem:local}. Thus $K$ is smaller than $G$ in the order of
Section~\ref{sec:red}: removing the leaves leaves $\beta$ unchanged
and reduces the edge count. Let $\delta=M(G)-M(K)$. Repeated leaf
deletion and minimality give
\begin{equation}\label{eq:delta}
 0\le\delta<\rho(G)-\rho(K)\le1 .
\end{equation}
The last bound holds because a center of $K$ reaches each added leaf
within $\rho(K)+1$. If $\Delta=3$, every core vertex has full degree
three. Its only possible lower-degree neighbor is its leaf, giving
$\ell_v=1$ at every core vertex. Equation~\eqref{eq:equal} gives a
zero change at every core edge, a contradiction. Hence $\Delta\ge4$.

\begin{proposition}[Degree-cut estimate]\label{prop:cut}
For an integer $3\le k<\Delta$ let $S_k=\{v\in V(K):d_v\le k\}$,
and let $N_k$ count edges of $K$ with exactly one end in $S_k$. Then
\begin{equation}\label{eq:cut}
 N_k\le k(k-1)\delta .
\end{equation}
\end{proposition}
\begin{proof}
Write $S=S_k$. Orient each unequal-degree core edge from its lower
endpoint to its higher endpoint. Denote by $I$ the directed edges
inside $S$ and by $F$ the directed edges leaving $S$. Every crossing
edge leaves $S$. Let $e_v$ count the core neighbors with degree $d_v$.
At $v$ there are $\ell_v-p_v$ incoming and $h_v-e_v$ outgoing edges.
Since $h_v a_v=\ell_v b_v$, we have
\[
 p_v b_v=(h_v-e_v)a_v-(\ell_v-p_v)b_v+e_v a_v .
\]
Sum this identity over $v\in S$. An internal directed edge $u\to v$
contributes $a_u-b_v$, and an edge leaving $S$ contributes $a_u$:
\begin{equation}\label{eq:flow}
 \sum_{v\in S}p_v b_v
 =\sum_{(u,v)\in I}(a_u-b_v)+\sum_{(u,v)\in F}a_u
        +\sum_{v\in S}e_v a_v .
\end{equation}
Each term of the first sum is positive by Lemma~\ref{lem:local}.

Delete precisely the leaves whose support belongs to $S$ and call the
result $G_S$. To compute the change of $M$, first consider a deleted
leaf at $v$. It has weight $1/d_v$, and the $\ell_v-1$ edges to
strictly lower core neighbors each gain $1/g_v$. Indeed, on a strict
edge with original degrees $d<e$, the new maximum remains $e$ when
the higher endpoint keeps its leaf or has none. If the higher endpoint
loses a leaf, the new maximum is $e-1$, as $d\le e-1$, irrespective
of any deletion at the lower endpoint. Thus the loss from that leaf,
after these gains, is $1/d_v-(\ell_v-1)/g_v=b_v$.

On a core edge of equal degrees $d$, the maximum decreases precisely
when both endpoints lose a leaf. In that case the edge gains
$1/[d(d-1)]$. Let $Q$ consist of these unordered core edges. The
exact total change is therefore
\begin{equation}\label{eq:subset}
 M(G)-M(G_S)=\sum_{v\in S}p_v b_v-\sum_{uv\in Q}\frac1{g_u}.
\end{equation}
For $uv\in Q$, the last sum of \eqref{eq:flow} counts $a_u+a_v$
on that edge. This is at least $2/g_u$ since $g_u=g_v$ and
$\ell_u,\ell_v\ge1$. All other contributions to that last sum are
nonnegative. Subtracting the correction in \eqref{eq:subset} from
\eqref{eq:flow}, and then deleting the remaining leaves one at a time
by Lemma~\ref{lem:delete}, gives
\[
 \sum_{(u,v)\in F}a_u\le M(G)-M(G_S)\le M(G)-M(K)=\delta .
\]
For every $(u,v)\in F$, one has $\ell_u\ge1$ and $d_u\le k$, so that
$a_u\ge1/[k(k-1)]$. There are $N_k$ such edges, giving \eqref{eq:cut}.
\end{proof}

\section{Summing the cuts and bounding the absolute radius}\label{sec:sum}
\begin{proof}[Completion of the proof of Theorem~\ref{thm:abs}]
Continue with the hypothetical minimal counterexample. Write
$t=\sum_{v\in V(K)}p_v$, so $G$ has $n+t$ vertices. In a graph
without isolated vertices, the sum of $1/d_u+1/d_v$ over the edges
counts each vertex once. Subtracting $2M(G)$ gives
\[
 n+t-2M(G)=\sum_{uv\in E(G)}\left|\frac1{d_u}-\frac1{d_v}\right|.
\]
A leaf at $v$ contributes $1-1/d_v$. A core edge with degrees
$d<e$ contributes $\sum_{k=d}^{e-1}1/[k(k+1)]$, counting exactly
the cuts it crosses. Equal-degree edges contribute zero. We obtain
\begin{equation}\label{eq:sum}
 n-2M(G)=\sum_{k=3}^{\Delta-1}\frac{N_k}{k(k+1)}
                     -\sum_{v\in V(K)}\frac{p_v}{d_v}.
\end{equation}
When all leaves are deleted, the retained core-edge weights can only
increase. Thus the sum of the removed weights $\sum_v p_v/d_v$ is at
least $\delta$. By Proposition~\ref{prop:cut},
\begin{align}
 n-2M(G)
 &\le \left(\sum_{k=3}^{\Delta-1}\frac{k-1}{k+1}-1\right)\delta
 \le(\Delta-4)\delta
 \le\Delta-4 . \label{eq:lowerM}
\end{align}
There are $\Delta-3$ summands, each at most one; the last two
inequalities use $\delta\ge0$, $\Delta\ge4$ and \eqref{eq:delta}.

Choose a core vertex of full degree $\Delta$. It has at least
$\Delta-1$ neighbors in $K$, since at most one is a leaf. By
Lemma~\ref{lem:tree}, $K$ has a spanning tree of diameter at most
$n-(\Delta-1)+1=n-\Delta+2$. Restore each original leaf and its
incident edge to that tree. Between any two vertices of the resulting
tree $T$, the path consists of a core-tree path and at most two leaf
edges. Thus $\diam(T)\le n-\Delta+4$. Distances in $G$ from a
diametral midpoint of $T$ are at most their distances in $T$. Applying
Lemma~\ref{lem:tree} and \eqref{eq:lowerM} gives
\begin{equation}\label{eq:contradiction}
 2\rho(G)\le\diam(T)\le n-\Delta+4\le 2M(G).
\end{equation}
This contradicts $M(G)<\rho(G)$ and proves Theorem~\ref{thm:abs}.
\end{proof}

\section{From absolute radius to vertex radius}\label{sec:finish}
Let $N=|V(G)|$, $r=\rad(G)$, and assume $\Delta\ge2$. We first keep
the difference $r-\rho(G)$ in the estimate; on graphs with a vertex
absolute center this term vanishes.
\begin{proposition}\label{prop:energy}
For each such graph,
\begin{equation}\label{eq:strong}
 R(G)-r \ \ge\
 \frac{N-2r-(\sqrt{\Delta}-1)(r-\rho(G))}{\sqrt{\Delta}+1}.
\end{equation}
\end{proposition}
\begin{proof}
The degree-square identity of Caporossi, Gutman, Hansen and
Pavlovi\'c \cite{CGHP} expresses $N-2R(G)$ as a sum of squares.
For clarity we derive it here and compare the summands, following the
bounded-degree argument in Section~6 of \cite{CPS}. On an edge with
degrees $d\le e$, the contribution to $R(G)-M(G)$ is
$1/\sqrt{de}-1/e\ge0$, and factoring gives
\[
 (d^{-1/2}-e^{-1/2})^2
  =(\sqrt{e/d}-1)\left(\frac1{\sqrt{de}}-\frac1e\right).
\]
Summing the expanded squares counts each vertex once in the terms
$1/d_u+1/d_v$. Since $\sqrt{e/d}\le\sqrt{\Delta}$, we have
\[
 N-2R(G)=\sum_{uv\in E(G)}(d_u^{-1/2}-d_v^{-1/2})^2
 \le(\sqrt{\Delta}-1)(R(G)-M(G)).
\]
Put $s=R(G)-r$ and $\eta=r-\rho(G)$. Theorem~\ref{thm:abs} gives
$R(G)-M(G)\le s+\eta$, whence
$N-2r-2s\le(\sqrt{\Delta}-1)(s+\eta)$. Solving for $s$
gives \eqref{eq:strong}.
\end{proof}

\begin{proof}[Proof of Theorem~\ref{thm:main}]
Apply Lemma~\ref{lem:tree} with $j=\Delta$ to obtain $N-2r\ge\Delta-2$.
Lemma~\ref{lem:half} gives $0\le r-\rho(G)\le1/2$.
Substitution in \eqref{eq:strong} yields
\[
 R(G)-r\ge
 \frac{\Delta-2-(\sqrt{\Delta}-1)/2}{\sqrt{\Delta}+1}
 =\sqrt{\Delta}-\tfrac32 .
\]
This is strictly positive for $\Delta\ge3$. A connected graph of
maximum degree two is a path or a cycle: following successive edges
either closes a cycle, which leaves no available edge to another
vertex, or ends at the two ends of a spanning path. For $N\ge3$ a
path has two edges of degrees $(1,2)$ and $N-3$ of degrees $(2,2)$.
In a cycle all edge degrees are $(2,2)$. Their vertex radii equal
$\lfloor N/2\rfloor$, by Lemma~\ref{lem:tree} on a path and by measuring
the two arc lengths on a cycle. The resulting values are
\[
\begin{array}{c|cc}
 & N\text{ even}&N\text{ odd}\\ \hline
 R(P_N)-\rad(P_N) & \sqrt2-3/2 & \sqrt2-1\\
 R(C_N)-\rad(C_N) & 0 & 1/2
\end{array}
\]
The singleton has $R=r=0$ and the single edge has $R=r=1$.
These cases exhaust all connected simple graphs. Since
$\sqrt2-3/2<0$, the first entry gives exactly the exceptional paths
and the global minimum.
\end{proof}

The coefficient of $\sqrt{\Delta}$ in \eqref{eq:main} cannot be
increased while retaining a constant additive term. For $\Delta\ge3$,
subdivide one edge of the star $K_{1,\Delta}$ once. Its diameter is
three and its vertex radius is two. There are $\Delta-1$ edges of
degrees $(1,\Delta)$, one of degrees $(2,\Delta)$ and one of degrees
$(1,2)$. For this family,
\[
 R(G)-r=\sqrt{\Delta}-2+\frac1{\sqrt2}
            -\frac{1-1/\sqrt2}{\sqrt{\Delta}} .
\]

\begin{thebibliography}{9}
\bibitem{CH}
G. Caporossi and P. Hansen.
Variable neighborhood search for extremal graphs 1: The AutoGraphiX system.
\emph{Discrete Mathematics} \textbf{212} (2000), 29--44.
\href{https://doi.org/10.1016/S0012-365X(99)00206-X}{doi:10.1016/S0012-365X(99)00206-X}.

\bibitem{CGHP}
G. Caporossi, I. Gutman, P. Hansen and L. Pavlovi\'c.
Graphs with maximum connectivity index.
\emph{Computational Biology and Chemistry} \textbf{27} (2003), 85--90.
\href{https://doi.org/10.1016/S0097-8485(02)00016-5}{doi:10.1016/S0097-8485(02)00016-5}.

\bibitem{CPS}
M. Cygan, M. Pilipczuk and R. \v{S}krekovski.
On the inequality between radius and Randi\'c index for graphs.
\emph{MATCH Communications in Mathematical and in Computer Chemistry}
\textbf{67} (2012), 451--466.
\href{https://match.pmf.kg.ac.rs/electronic_versions/Match67/n2/match67n2_451-466.pdf}{Journal PDF}.

\bibitem{Deng}
H. Deng, Z. Tang and J. Zhang.
On a conjecture of Randi\'c index and graph radius.
\emph{Filomat} \textbf{29} (2015), 1369--1375.
\href{https://doi.org/10.2298/FIL1506369D}{doi:10.2298/FIL1506369D}.

\bibitem{Doig}
M. I. Doig.
Block eccentricity, a radius bound, and an application to the Randi\'c index.
Preprint (2023). \href{https://arxiv.org/abs/2309.11613}{arXiv:2309.11613}.

\bibitem{DLS}
Z. Dvo\v{r}\'ak, B. Lidick\'y and R. \v{S}krekovski.
Randi\'c index and the diameter of a graph.
\emph{European Journal of Combinatorics} \textbf{32} (2011), 434--442.
\href{https://doi.org/10.1016/j.ejc.2010.12.002}{doi:10.1016/j.ejc.2010.12.002}.

\bibitem{Faj}
S. Fajtlowicz.
On conjectures of Graffiti.
\emph{Discrete Mathematics} \textbf{72} (1988), 113--118.
\href{https://doi.org/10.1016/0012-365X(88)90199-9}{doi:10.1016/0012-365X(88)90199-9}.

\bibitem{VMRC}
Keio.
\emph{Verifiable Mathematical Reasoning Corpus}.
Research drafts [087], [094] and [106].
\href{https://github.com/QEEDDD/Verifiable-Mathematical-Reasoning-Corpus/tree/8a3f6895d0b9b9f125e43157e4a003a2a2b51704}{Repository snapshot 8a3f6895d0b9}.

\bibitem{Zuo}
L. Zuo.
About a conjecture on the Randi\'c index of graphs.
\emph{Bulletin of the Malaysian Mathematical Sciences Society}
(2) \textbf{35} (2012), 411--424.
\href{https://math.usm.my/bulletin/pdf/v35n2/v35n2p16.pdf}{Journal PDF}.
\end{thebibliography}
\end{document}

